\section{Introduction}

%From the first perturbative expansions in quantum systems
%with finite degrees of freedom, to its modern formulation in QFT, perturbation
%theory forces us to question the foundamental belief that, to understand the dynamical world
%that surrounds us, we need to ground its description in the simplest mechanisms available, and
%try to recover the ``macroscopic" phenomena as limits of a more elementary theory.

The rigorous analysis of many body systems poses notoriously hard challanges that push us to
develop more and more sofisticated tools, and to ask ourselves the true meaning of the
mathematical objects that come out in the process. One of the most important and useful tools
in theoretical physics that allows for an extremely close connection between theory and
experiment is perturbation theory. It allows -- among other things -- for the analysis of
coarse properties of a system starting from its most foundamental interactions, and his comes
as an especially powerful process considering the mathematics behind quantum mechanics seems
to get more and more complicated and intricate the deeper we analyze a system. Even just
tackling the conceptually simple case of a gas with pair-wise interactions proves to be
mathematically extremely challenging, and it is precisely this difficulty that motivates our
work.

For a system of many interacting particles it can become a nightmare to discuss the
effect of every interaction a single particle may have with the rest of the system, and it is
therefore a good idea to first try looking for meachanisms that either make the system easier
to study or the mathematics better-behaved. An example that is very easy to picture is the
screening of the Coulomb interaction between the nucleus of an atom and its orbiting electrons,
be it in single atoms, molecules, crystal lattices, and so on.
It is not at all a trivial endevour trying to compute exactly even just the ground state
energy of the atom, but we can make some useful observations before attempting to do so. If we
consider electrons in the innermost shells we can expect their binding energy to be mainly due
to the strong electrostatic interaction with the protons inside the nucleus. But if we move to
higher and higher energy levels, the electrons, while still being attracted by the nucleus,
will also be repelled by the presence of more and more lower-energy electrons, which effectively
\textit{screen} the total charge of the nucleus \cite[p.~51]{ManBook}. Thus, electrons in higher
shells will not
feel the full long-range effect of the Coulomb interaction, but rather an effective
interaction that is suppressed depending on their energy level, assuming lower energy levels
are occupied.
\\
The problem then becomes how this screened interaction behaves. This is where \textit{mean
field methods} come into play. The easy example is once again that of many-electron atoms. To
find the ground state energy without directly accounting for the very complicated correlations
between each particle, one starts by assuming that electrons are found in independent
single-electron atomic
orbitals. Then, computing the overall effect of such a configuration on the energy of one of
the electrons, one can look for a better configuration that minimizes this effect, and
iterating this procedure leads to the Hartree-Fock approximation for electronic configurations
\cite[pp.~48--50]{ManBook}. Hartree-Fock theory can be formulated more generally and applied
to much more complicated settings. Furthermorre, it essentially corresponds to the first order
contribution to the ground state energy in the case of homogeneous systems \cite[p.~81]{GiuVigBook}.

This screening phenomenon is, of course, not limited to electrons in atomic orbitals. In any
system with many interacting particles, one may ask whether the state of each particle can be
described not by the true, \textit{bare}, interaction that characterizes the system, but
rather by a rescaled, \textit{dressed}, interaction. If this is the case, it may also happen
that ill-behaved quantities become well-behaved. As an example, in the case of an electron gas
it is precisely the long range of the Coulomb interaction that makes perturbation theory badly
divergent in the IR \cite[p.~38--39]{GiuVigBook}.
In the computation of the ground state energy, the RPA comes as a way of carrying out such
a computation once the screening of the interaction due to the high number of
particles is accounted for; indeed, today we also understand the RPA as a partial resummation
of perturbation theory at all orders. This resummation, coincidentally, also removes the
divergence at low momentum of the Coulomb interaction \cite[p.~293]{GiuVigBook}).

These are just some examples of how perturbation theory can help us see how our physical
understanding of the simpler mechanisms should be gauged by our knowledge of the range of
applicability of a theory. Ill-defined quantities, computations that are too difficult to carry
out, effective phenomena that present themselves differently depending on energy scales, they
can all be understood through the lenses of perturbation theory as foundamental tool. But
perturbation theory alone is not enough: it needs to be carefully controlled and, at some
point, perturbative results must be rigorously proved to be fully validated.
\textcolor{red}{Alla faccia dei fisici delle particelle.}

We now set out to introduce the mathematical formalism necessary to understand perturbation
theory, and use it to prove -- perturbatively, of course -- that the mean field scaling
limit we take in exam does reproduce, at second order, the known exact results for the ground
state energy of an interacting Fermi gas. Our hope is that, by laying the groundwork with the
present results, the full perturbative series can be treated and further developments can be
achieved through renormalization group methods to tackle the bigger problem of interacting
fermions.
